\chapter{Code Execution Sandbox}

\section{Purpose and Design}
Coding questions are only useful if the learner can check an answer. \texttt{recall/sandbox.py} provides this by running submitted code against the test cases that accompany each coding question. Each test case consists of a stdin string and the expected stdout.

\section{Execution Model}
For every test case the module performs the following steps:
\begin{enumerate}
\item The code is written to a temporary \texttt{.py} file.
\item It is executed with the same Python interpreter as the server (\texttt{sys.executable}) in a subprocess, with the test input piped to stdin and the environment variable \texttt{PYTHONIOENCODING=utf-8}.
\item Standard output and standard error are captured and truncated to 10{,}000 characters each.
\item The subprocess is limited to a 5-second time-out; on expiry it is terminated and the case is reported as timed out.
\item The temporary file is always removed in a \texttt{finally} block.
\end{enumerate}

\section{Result Evaluation}
A test case passes only if the process did not time out, exited with return code~0 and its trimmed stdout equals the trimmed expected output. For each case the endpoint returns the input, the expected and actual output, a pass flag and an error message (a time-out notice or the last error text from stderr). The response also carries the totals \texttt{passed}, \texttt{failed} and \texttt{total}.

\begin{lstlisting}[style=pseudo, caption={Evaluation of one test case}, label={lst:tc}]
result <- run_code(code, stdin = test.input, timeout = 5s)
actual <- strip(result.stdout)
passed <- (not result.timed_out) and result.returncode == 0
          and actual == strip(test.expected_output)
\end{lstlisting}

\section{Language Support and Limitations}
The user interface lets the learner request coding questions in Python, Java, C, C++ or JavaScript, but \textbf{only Python can currently be executed}; for any other language the runner returns the message ``Only Python execution is supported currently''. The runner is a convenience for local practice and has the following security limits, which are discussed further in Chapter~\ref{chap:security}:
\begin{itemize}
\item the code runs with the same privileges as the server, without file-system or network isolation;
\item only wall-clock time and output size are limited, not memory or CPU.
\end{itemize}
It must therefore never be exposed to untrusted users on a public network.

\screenshot{coding_question}{A coding question with starter code editor and test-case results after pressing Run}{fig:coding}{Generate a quiz with one Python coding question, type or paste a solution, press Run and capture the editor plus the passed/failed test list.}
